% 20260708(3)_Assingment_ComputationalMathematics_A
\documentclass[dvipdfmx]{jsreport}
%preamble
\usepackage{soupreamble}

\title{計算数学特論A 第7回レポート課題}
\author{M1 6012 川村聡一郎}
\date{}


\begin{document}
\maketitle

\begin{tcolorbox}
\begin{enumerate}
	\item 独立集合(判定版)から最大クリーク(判定版)への多項式時間還元が存在することを示せ，
	\item Hamiltonパス問題(判定版)から最長単純パス問題(判定版)への多項式時間還元が存在することを示せ，
\end{enumerate}
\end{tcolorbox}

\begin{souanswer} %答案はここに書く
\begin{enumerate}
	\item 独立集合(判定版)の入力であるグラフを $G=(V, E)$，自然数を $k$ とする，
	次のような最大クリーク問題の入力 $(G', k')$ を構成する，
\begin{itemize}
	\item $G'= (V, E')$\ \ ただし$E'$は$G$の補グラフ
	\item $k'= k$
\end{itemize}
	$S\subset V$が$G$におけるサイズ$k$以上の独立集合であると仮定する，独立集合の定義より任意の任意の $u, v \in S$ について $\{u, v\} \notin E$ である，
	これは，補グラフ $G'$ の定義より，任意の $u, v \in S$ について $\{u, v\} \in E'$ であることと同値である，\\
	すなわち，$S$ が $G$ の独立集合であることと，$S$ が $G'$ のクリーク（完全部分グラフ）であることは同値である，\\
	ここで$k' = k$ と設定しているため，$G$ にサイズ $k$ 以上の独立集合が存在することと，$G'$ にサイズ $k'$ 以上のクリークが存在することは同値となり，還元の正当性が示された，  【多項式時間の確認】補グラフ $G'$ の頂点集合は元の $V$ と同一であり，辺集合 $E'$ の構成はすべての頂点ペア（$O(|V|^2)$ 個）について隣接関係を反転させるだけで構築できる，したがって，入力グラフ $G$ のサイズの多項式時間で還元が可能である，

	\item ハミルトンパス問題の入力であるグラフを $G=(V, E)$ とする，
	次のような最長単純パス問題の入力 $(G', C)$ を構成する，
\begin{itemize}
	\item $G'= C$
	\item $C= |V|- 1$
\end{itemize}
	定義より，長さ$k$の単純パスは$k+ 1$個の相異なる頂点の列からなる，したがって，長さ$C$の単純パスは$C+ 1= (|V|- 1)+ 1= |V|$個の相異なる頂点を通るグラフ$G'$に長さ$C$以上の単純パスが存在すると仮定する，グラフの総頂点数は $|V|$ 個しかないため，このパスはちょうど $|V|$ 個の頂点をすべて1回ずつ通るパスとなる，\\
	ハミルトンパスは「全ての頂点を通る単純パス」と定義されているため，これはハミルトンパスそのものである，\\
	逆に，$G$ にハミルトンパスが存在すれば，それは $|V|$ 個の相異なる頂点を通るパスであり，その長さは $|V|- 1= C$ となるため，$G'$ に長さ $C$ 以上の単純パスが存在する，この還元において，グラフ $G'$ は入力グラフ $G$ をそのまま出力するのみである，\\
	また，出力する自然数$C$の値は$|V|- 1$であるが，最長単純パス問題への「自然数のサイズ」はそのビット長で評価される$C$を2進数で表現したときのサイズは$O(\log|V|)$となり，入力グラフの「グラフのサイズ(頂点数$|V|$や辺数$|E|$)」に対して対数オーダーに収まる，\\
	グラフのコピーおよび頂点数のカウントと$C$の算出は線形時間で実行可能であり，出力されるデータサイズも入力サイズの多項式に収まるため，この還元は正しく多項式時間還元となっている，
\end{enumerate}
\end{souanswer}












\end{document}